\section{The French Hermetic Tradition}

Valentin Tomberg explains why he wrote his meditations on the major arcana of the Tarot in French:

\begin{quotex}
These letters were written in French, which is not the mother tongue of the author, because it is in France, and in France only, that a living literature on the Tarot has been perpetuated since the 18th century. Furthermore, there exists as well a continuous tradition of Hermetism, which unites a spirit of free research with respect for the Tradition. These letters, by virtue of their contents, will therefore be able to become an integral part of the Tradition while enriching it.\footnote{My translation.} 
\end{quotex}


Some of the key figures in this tradition are mentioned here.

\begin{quotex}
I also have some Known Friends, but most of them are found in the spiritual world. That is the main reason that I address myself to them in these letters. And how many times, while writing them, have I felt the fraternal embrace of these friends which include Papus\footnote{\url{http://en.wikipedia.org/wiki/G\%C3\%A9rard\_Encausse}}, Guaita\footnote{\url{http://en.wikipedia.org/wiki/Stanislas\_de\_Guaita}}, Peladan\footnote{\url{http://en.wikipedia.org/wiki/Peladan}}, Eliphas Levi\footnote{\url{http://en.wikipedia.org/wiki/Eliphas\_Levi}}, and Claude de Saint-Martin\footnote{\url{http://en.wikipedia.org/wiki/Louis\_Claude\_de\_Saint-Martin}}!\footnote{My translation.}
\end{quotex}


Following Claude de Saint-Martin, Peladan, and Eliphas Levi, Tomberg affiliated himself with the Catholic Church, for reasons that will be explained in a subsequent post. Tomberg, however, knew none of these men personally.

Rene Guenon, on the other hand, did know Papus, and probably Peladan and Guaita personally, since he was active, in his youth, in the French Hermetic revival initiated by those men. However, he broke with Papus, ostensibly over the doctrine of reincarnation. Subsequently, he never came to terms with Hermetism and did not regard it as a full tradition, but rather as mere cosmology. This led him to dismiss Evola's work on the Hermetic Tradition.

For a time, it seems Guenon was attempting to influence Catholicism in a more metaphysical direction, since he regarded the Medieval religion as an authentic Tradition. However, he was rebuffed, primarily by the French theologian Jacques Maritain. So here we have two esoterists: the Russian working within the French, Catholic, Hermetic Tradition, and the Frenchman rejecting it in favour of Sufism.

An intermediate view is taken by Julius Evola in his \emph{The Hermetic Tradition}. While joining with Guenon in his understanding of metaphysics, he embraces the Hermetic Tradition both as a means of practical work and as a the Tradition of the West. Also, unlike Tomberg, Evola sees the Hermetic Tradition as transcending the Christian symbolism of many of the Hermetists. Nevertheless, he refers to Eliphas Levi to summarize the purpose of Hermetism.

\begin{quotex}
Eliphas Levi repeatedly warns that we must make use of every hour and moment; we must rid the will of any dependency and accustom it to domination, we must become absolute master of the self, learn how to deny that call of pleasure, hunger, and sleep and be unmoved by success or failure. Life must be given over to the will directed by one's thought and served by the entire nature \emph{to submit in every organ to the spirit}, and by sympathy to all universal forces corresponding to them. All the faculties and senses must participate in the Work; nothing must be left inactive. The genuine spirit must have secured itself against all danger of hallucination and terror, and so find us purified inwardly and outwardly.
\end{quotex}



\flrightit{Posted on 2010-09-13 by Cologero}

\begin{center}* * *\end{center}

\begin{footnotesize}
\begin{sffamily}
\texttt{Elsinore on 2010-09-14 at 00:25 said:}

I would have included Fulcanelli in this article as well. Granted, we may never know his true nationality but his language of choice for the transmission of his lessons was French.

\hfill

\texttt{Cologero on 2010-09-14 at 06:54 said:}

I was quoting from Tomberg, whose personal list was more narrow. But certainly, we could add Fulcanelli, and Schwaller de Lubicz as well. And I would add Joseph de Maistre.

\hfill

\texttt{Cologero on 2010-09-14 at 18:33 said:}

Then there is Maitre Philippe, often mentioned by Tomberg. He is interesting especially since he spent some time in St Petersburg, Tomberg's birthplace. St Petersburg had a Martinist presence, thus a ``French connection''. Goes back to Maistre, a Martinist, who also spent many years there.

Two other important writers, both for Tomberg and Guenon, are Fabre d'Olivet and Saint-Yves d'Alveydre. The latter's writing on Authority and Power were adopted by both Tomberg and Guenon. Another writer is Matgioi, often mentioned by Guenon. Matgioi was a student of Chinese spirituality, and gave Guenon the term ``Primordial Tradition''. (Although the Catholic Joseph de Maistre had a similar notion.)


\hfill

\end{sffamily}
\end{footnotesize}

